\section{Conclusion and Future Work}
\label{sec:conclusion}

This paper introduced \toolname{}, a novel framework that combines large language models with Verus to automatically verify Rust data-structure modules. \toolname{} systematically generates \texttt{View} implementations, type invariants, specifications, and proof code under the scheduling of an extra planner module. To mitigate the issue that LLMs do not fully grasp Verus' annotation syntax and verification-specific semantics, we not only inject syntax guidelines in the prompt, but also introduce an additional repair stage to fix errors.%
\ifarxiv
\ Table~\ref{tab:repair-heuristics} in the appendix details the specialized heuristics that support these repairs.
\fi

We evaluated \toolname{} on eleven benchmarks drawn from the Verus GitHub repository. Across eleven benchmark modules, \toolname{} successfully verifies ten of them and verifies 128/129 (99.2\%) of functions in our benchmark. \toolname{} shows a significant improvement over the baseline, demonstrating the framework's effectiveness.



% while
%substantially reducing the handwritten annotations required compared to
%manual development. The resulting artifacts demonstrate that
%LLM-guided synthesis can substantially lower the effort required to construct
%formally verified data structures in practice.


%By decomposing the construction of proofs into generation and repair stages, the system . Furthermore, we inject syntax guidelines in the 

%This planner-guided process turns raw prover
%errors into actionable repair steps, enabling the framework to converge on
%verified implementations with limited human input; Table~\ref{tab:repair-heuristics}
%details the specialized heuristics that support these repairs.


%Looking ahead, extending the framework with richer lemma retrieval, support for
%concurrent resource algebras, and automatic unit-test generation could further
%increase its automation power. These enhancements would broaden the range of
%Rust libraries amenable to verification and bring mechanically checked
%correctness guarantees within reach for an even wider set of developers.
